% Part: incompleteness
% Chapter: incompleteness-provability
% Section: second-incompleteness-thm

\documentclass[../../../include/open-logic-section]{subfiles}

\begin{document}

\olfileid{inc}{inp}{2in}

\olsection{نامکملیت دا دوجا قضیہ}

اسیں ایہہ گل کیویں بیان کریے کہ $\Th{PA}$ اپنی بے تضادی نوں
ثابت نہیں کردا؟ $\Th{PA}$ نوں متضاد کہنا ایہہ کہنا اے کہ
$\Th{PA} \Proves \eq[0][1]$۔ ایس لئی بے تضادی دے جملے
$\OCon[\Th{PA}]$ نوں !!{sentence} $\lnot
\OProv[\Th{PA}](\gn{\eq[0][1]})$ دے روپ وچ لے سکدے آں؛فیر
ہِٹھلا قضیہ مطلوبہ نتیجہ دیندا اے:

\begin{thm}
\ollabel{thm:second-incompleteness} 
جے $\Th{PA}$ بے تضاد اے تاں $\Th{PA}$، $\OCon[\Th{PA}]$
نوں !!{derive} نہیں کردا۔
\end{thm}

ایہہ دھیان وچ رکھنا اہم اے کہ قضیہ $\OCon[\Th{PA}]$ دی خاص
نمائندگی اُتے منحصر اے،یعنی $\OProv[\Th{PA}](y)$ دی خاص
نمائندگی اُتے۔ سانوں صرف ایہہ چاہیدا اے کہ
$\OProv[\Th{PA}](y)$ دی نمائندگی تین !!{derivability} شرطاں
پوری کرے؛ایس لئی قضیہ ہر اُس نظریے تک وی پھیلدا اے جس دا
!!a{derivability} محمول ایہناں خاصیتاں نوں پورا کردا ہووے۔

گوڈل دی دلیل دا خاکہ پڑھنا فائدہ مند اے،کیونکہ قضیہ چنگے
اختتامی نکتے وانگ سامنے آ جاندا اے۔ دلیل ایویں اے۔
$!G_\Th{PA}$ اوہ گوڈل جملہ ہووے جیہڑا اسیں
\olref[1in]{thm:first-incompleteness} دے ثبوت وچ بنایا سی۔ اسیں
دکھا چکے آں: «جے $\Th{PA}$ بے تضاد اے تاں $\Th{PA}$،
$!G_\Th{PA}$ نوں !!{derive} نہیں کردا»۔ جے اس گل نوں
$\Th{PA}$ دے \emph{اندر} رسمی بنائیے،تاں ہِٹھلا نتیجہ ثابت
ہوندا اے:
\[
\OCon[\Th{PA}] \lif \lnot \OProv[\Th{PA}](\gn{!G_\Th{PA}}).
\]
% PNB-CMP-071: B9 bare Prov دی تھاں formal OProv؛اصل اگلی تمام صورتاں نال مطابقت۔
ہن منّو $\Th{PA}$، $\OCon[\Th{PA}]$ نوں !!{derive}s کردا
اے۔ فیر اوہ $\lnot
\OProv[\Th{PA}](\gn{!G_\Th{PA}})$ نوں وی !!{derive}s کردا اے۔
پر کیوں جے $!G_\Th{PA}$ گوڈل جملہ اے،ایہہ
$!G_\Th{PA}$ نال ہم ارز اے۔ ایس لئی
$\Th{PA}$، $!G_\Th{PA}$ نوں !!{derive}s کردا اے۔

پر اسیں جاندے آں کہ جے $\Th{PA}$ بے تضاد ہووے تاں اوہ
$!G_\Th{PA}$ نوں !!{derive} نہیں کردا!{} سو جے $\Th{PA}$
بے تضاد اے تاں اوہ $\OCon[\Th{PA}]$ نوں !!{derive} نہیں
کر سکدا۔

دلیل نوں ہور دقیق بناؤن لئی $!G_\Th{PA}$ نوں
$\Th{PA}$ لئی گوڈل جملہ منّاں گے،تے !!{derivability} شرطاں
(P1)--(P3) نال دکھاواں گے کہ $\Th{PA}$،
$\OCon[\Th{PA}] \lif
!G_\Th{PA}$ نوں !!{derive}s کردا اے۔ ایس توں پتا لگے گا کہ
$\Th{PA}$، $\OCon[\Th{PA}]$ نوں !!{derive} نہیں کردا۔
ہِٹھاں $\Th{PA}$ دے اندر ایس ثبوت دا خاکہ اے۔ (سہولت لئی
اسیں $\Th{PA}$ دے ذیلی اشاریے چھڈ دیندے آں۔)
\begin{align}
& !G \liff \lnot \OProv(\gn{!G}) \ollabel{G2-1}\\
& \qquad\text{$!G$، گوڈل دا جملہ اے}\notag \\
& !G \lif \lnot \OProv(\gn{!G}) \ollabel{G2-2}\\
  & \qquad\text{\olref{G2-1} توں} \notag\\
& !G \lif
  (\OProv(\gn{!G}) \lif \lfalse) \ollabel{G2-3}\\
  & \qquad\text{\olref{G2-2} توں منطق نال}\notag\\
& \OProv(\gn{
    !G \lif
    (\OProv(\gn{!G}) \lif \lfalse)
  }) \ollabel{G2-4}\\
  & \qquad\text{\olref{G2-3} توں شرط P1 نال} \notag\\
& \OProv(\gn{!G}) \lif
  \OProv(\gn{
    (\OProv(\gn{!G}) \lif \lfalse)
    }) \ollabel{G2-5}\\
  & \qquad\text{\olref{G2-4} توں شرط P2 نال} \notag\\
& \OProv(\gn{!G}) \lif (\OProv(\gn{\OProv(\gn{!G})}) \lif \OProv(\gn{\lfalse})) \ollabel{G2-6}\\
  & \qquad\text{\olref{G2-5} توں شرط P2 تے منطق نال} \notag\\
& \OProv(\gn{!G}) \lif 
  \OProv(\gn{\OProv(\gn{!G})}) \ollabel{G2-7}\\
   & \qquad\text{شرط P3 نال} \notag\\
& \OProv(\gn{!G}) \lif \OProv(\gn{\lfalse}) \ollabel{G2-8}\\
  & \qquad \text{\olref{G2-6} تے \olref{G2-7} توں منطق نال}\notag\\
& \OCon \lif \lnot \OProv(\gn{!G}) \ollabel{G2-9}\\
  & \qquad\text{\olref{G2-8} دے عکسِ نقیض تے $\OCon \ident \lnot \OProv(\gn{\lfalse})$ نال}\notag \\
& \OCon \lif !G \notag\\
  & \qquad\text{\olref{G2-1} تے \olref{G2-9} توں منطق نال}\notag
\end{align}
اُتلی دلیل وچ منطق دی ورتوں صرف قضیاتی منطق دے بنیادی حقائق
نیں۔ مثال لئی \olref{G2-3}،
$\Proves \lnot!A \liff (!A\lif
\lfalse)$ ورتدا اے،تے \olref{G2-8}،
$!A \lif (!B \lif !C), !A \lif !B
\Proves !A \lif !C$ ورتدا اے۔ \olref{G2-5} تے
\olref{G2-6} وچ شرط~P2 دی ورتوں P2 دیاں انہاں مثالاں اُتے
منحصر اے:
$\OProv(\gn{!A \lif !B}) \lif
(\OProv(\gn{!A}) \lif \OProv(\gn{!B}))$۔ پہلی مثال وچ
$!A \ident
!G$ تے $!B \ident \OProv(\gn{!G}) \lif \lfalse$؛
% PNB-CMP-071: B11 آخری P2 instance وچ G دے formula marker ! بحال۔
دوجی وچ $!A
\ident \OProv(\gn{!G})$ تے $!B \ident \lfalse$۔

نامکملیت دے دوجے قضیے دا ہور عمومی روپ ایہ اے:

\begin{thm}
\ollabel{thm:second-incompleteness-gen} منّو $\Th{T}$ کوئی وی
بے تضاد، !!{axiomatized} نظریہ اے جیہڑا $\Th{Q}$ دی توسیع
کردا اے؛تے $\OProv[\Th{T}](y)$ کوئی وی ایسا فارمولا ہووے
جیہڑا $\Th{T}$ لئی !!{derivability} شرطاں P1--P3 پوری کردا
اے۔ فیر $\Th{T}$، $\OCon[T]$ نوں !!{derive} نہیں کردا۔
\end{thm}

\begin{prob}
دکھاؤ کہ $\Th{PA}$، $!G_{\Th{PA}} \lif \OCon[\Th{PA}]$
نوں !!{derive}s کردا اے۔
\end{prob}

\begin{digress}
ایس گل دا حاصل ایہہ اے کہ ریاضی دا کوئی وی «معقول» بے تضاد
نظریہ اپنی بے تضادی دا جملہ !!{derive} نہیں کر سکدا۔ منّو
$\Th{T}$ ریاضی دا نظریہ اے جیہڑے وچ $\Th{Q}$ تے ہلبرٹ
دی «محدود وسیلی» دلیلاں (اوہناں توں جو وی مراد ہووے) شامل نیں۔
تاں پورا $\Th{T}$، $\Th{T}$ دی بے تضادی دا جملہ
!!{derive} نہیں کر سکدا؛سو لازماً اوہدا محدود
وسیلی حصہ وی $\Th{T}$ دی بے تضادی دا جملہ !!{derive} نہیں
کر سکدا۔ ایس معنی وچ «پوری ریاضی» لئی محدود وسیلی بے تضادی
دا ثبوت نہیں ہو سکدا۔

«محدود وسیلی» لفظ دی تعبیر وچ کجھ گنجائش اے،تے گوڈل
1931 دے مقالے وچ ایہہ امکان منّدا اے کہ جیہڑی گل نوں اسیں
«محدود وسیلی» سمجھدے آں اوہ ریاضی دیاں انہاں قسماں توں باہر
ہووے جنہاں نوں ہلبرٹ رسمی بنانا چاہندا سی۔ پر ایتھے گوڈل
نے کھُل دلی دکھائی؛اج ایہہ سمجھنا اوکھا اے کہ کوئی شے معقول
طور تے «محدود وسیلی» وی کہلائے تے،مثلاً $\Th{ZFC}$،یعنی
انتخاب دے مسلّمے سمیت زرمیلو--فرینکل سیٹ نظریے،وچ رسمی طور
بیان وی نہ کیتی جا سکے۔
\end{digress}

\end{document}
